\begin{frame}{Локальная оптимизация методом Нелдера--Мида}
	\vfill
	
	\begin{columns}[T]
		
		% --- ЛЕВАЯ КОЛОНКА: МАТЕМАТИЧЕСКАЯ РОЛЬ ---
		\begin{column}{0.49\textwidth}
			\textbf{Назначение и свойства метода:}
			\begin{tcolorbox}[
				colframe=DeepBlue!35, 
				colback=DeepBlue!1!Ivory, 
				coltext=DeepBlue, 
				boxrule=0.5pt, arc=3pt, 
				top=4pt, bottom=4pt, left=4pt, right=4pt,
				before=\vspace{3pt}, after=\vspace{0pt},
				equal height group=neldermeadstrict
				]
				\footnotesize
				\begin{itemize}
					\setlength{\itemsep}{5pt}
					\item Финишная локальная минимизация в окрестности найденного глобального приближения $\theta_{\text{PSO}}^*$.
					\item \textbf{Класс метода:} прямой метод безусловной оптимизации (алгоритм деформируемого многогранника).
					\item Быстрая сходимость в гладких локальных окрестностях без необходимости вычисления градиентов функционала $\mathcal{F}(\theta)$.
				\end{itemize}
			\end{tcolorbox}
		\end{column}
		
		% --- ПРАВАЯ КОЛОНКА: ПАРАМЕТРЫ ПОИСКА ---
		\begin{column}{0.49\textwidth}
			\textbf{Параметры численной реализации:}
			\begin{tcolorbox}[
				colframe=Burgundy!35, 
				colback=Burgundy!1!Ivory, 
				coltext=DeepBlue, 
				boxrule=0.5pt, arc=3pt, 
				top=4pt, bottom=4pt, left=4pt, right=4pt,
				before=\vspace{3pt}, after=\vspace{0pt},
				equal height group=neldermeadstrict
				]
				\footnotesize
				\begin{itemize}
					\setlength{\itemsep}{5pt}
					\item Начальный симплекс формируется в окрестности $\theta_{\text{PSO}}^*$ и включает $4+1 = 5$ вершин.
					\item Итерационное изменение формы симплекса через операторы отражения, растяжения и сжатия.
					\item Критерий останова — достижение сходимости по функционалу с точностью $\varepsilon = 10^{-7}$.
					\item \texttt{scipy.optimize}.
				\end{itemize}
			\end{tcolorbox}
		\end{column}
		
	\end{columns}
	
	\vfill
\end{frame}